\documentclass[conference]{IEEEtran}
\usepackage{cite}
\usepackage{amsmath,amssymb}
\usepackage{graphicx}
\usepackage{booktabs}
\usepackage{algorithm}
\usepackage{algpseudocode}

\graphicspath{{figures/}}

\begin{document}

\title{Sequential Causal Training of Windowed PINNs on the Lorenz-1960 Three-Wave System}

% \author{
% \IEEEauthorblockN{Ikramul Hasan Moral, Md. Abu Bakar, Samiur Rahman Omlan,\\
% Md. Touhidul Islam, Fariha Islam, and Muhammad Nomani Kabir}
% \IEEEauthorblockA{Department of Computer Science and Engineering, United International University, Dhaka, Bangladesh\\
% imoral223489@bscse.uiu.ac.bd, mbakar223200@bscse.uiu.ac.bd, somlan223195@bscse.uiu.ac.bd,\\
% mislam223435@bscse.uiu.ac.bd, fislam223431@bscse.uiu.ac.bd, kabir@cse.uiu.ac.bd}
% }


\author{\IEEEauthorblockN{Ikramul Hasan Moral}
\IEEEauthorblockA{\textit{Computer Science And Engineering} \\
imoral223489@bscse.uiu.ac.bd\\
0112230489}\\

\and
\IEEEauthorblockN{Md. Abu Bakar}
\IEEEauthorblockA{\textit{Computer Science And Engineering} \\
mbakar223200@bscse.uiu.ac.bd\\
0112230200}
\and
\IEEEauthorblockN{Samiur Rahman Omlan}
\IEEEauthorblockA{\textit{Computer Science And Engineering} \\
somlan223195@bscse.uiu.ac.bd\\
0112230195}
\and
\IEEEauthorblockN{Md. Touhidul Islam}
\IEEEauthorblockA{\textit{Computer Science And Engineering} \\
mislam223435@bscse.uiu.ac.bd \\0112230435}
\and
\IEEEauthorblockN{Fariha Islam}
\IEEEauthorblockA{\textit{Computer Science And Engineering} \\
fislam223431@bscse.uiu.ac.bd \\ 0112230431}
\and
\IEEEauthorblockN{Dr. Muhammad Nomani Kabir}
\IEEEauthorblockA{\textit{Professor, Dept of CSE} \\
\textit{United International University}\\
kabir@cse.uiu.ac.bd}
}


\maketitle

\begin{abstract}
We study a physics-informed neural network (PINN) approximation of a single full-cycle orbit of the Lorenz three-state system on $t\in[0,13.26446]$. We trained a single network on the full interval that produced a combined state root-mean-square error (RMSE) of 2.237. To address this long-horizon problem, we divide the interval into 27 windows. Each window has a separate network and an initial state that is fixed to the predicted endpoint of the previous window. In each window, Adam reduces the residual in four stages. In each stage, we use causal weights and causal weighted loss to reduce the loss of all earlier points. Then, L-BFGS minimizes the loss using the unweighted residual mean. A numerical reference is used to monitor the training but is not used for optimization or stopping. In a single-seed run, the windowed model gives a combined RMSE of $6.97\times10^{-5}$ on 13,265 evaluation times and a dense-grid residual mean square of $1.87\times10^{-7}$.  The full-interval and windowed runs also differ in parameter count, collocation sampling, and optimizer schedule; the comparison is therefore descriptive and does not isolate the effect of causal weighting by itself.

\end{abstract}

\begin{IEEEkeywords}
physics-informed neural network, Lorenz-1960 system, causal training, time windows, ordinary differential equation, windowed PINN
\end{IEEEkeywords}

\section{Introduction}
The Lorenz-1960 model is a three-state, nonlinear ordinary differential equation (ODE) obtained from a reduced flow model \cite{lorenz1960}. It has an initial-value problem which provides a simple test of the convergence of a learned continuous-time function from a specified trajectory over a long interval. A reference trajectory is provided by numerical integration. A PINN is trained with the residual of the ODE instead of the solution samples \cite{raissi2019}. Training a long-horizon problem with a traditional PINN is not recommended by previous works~\cite{wang2024causal,krishnapriyan2021}. It couples early and late times, but the state at a later time is determined by the previous trajectory. To overcome this ordering, Wang \emph{et al.} suggested the use of time-dependent residual weights PINN training \cite{wang2024causal}. An initial condition can also be met by construction using a hard trial solution  \cite{lagaris1998}. We use both concepts, we split the time period in equal 27 windows and a separate network for every windows. The initial state is fixed with the endpoint predicted in one window of the next. The parameters of a trained window is copied to the next trainable window. This paper provides an implementation and its saved results. Our contributions are threefold: (i) a config-locked, resumable training pipeline whose numerical reference is verified by conserved-quadratic and cross-solver checks; (ii) a quantitative characterization of two failure modes on the full-cycle interval, residual plateauing and handoff forgetting; and (iii) a descriptive comparison of the full-interval and 27-window causal runs, with a budget-matched causal ablation and a multi-seed error budget released with the run artifacts. An interval that completes a full-cycle loop for the chosen initial state is $[0,13.26446]$; we report the saved result on this essentially closed orbit. The actual training goal and window handoff is then given and compared to a saved single-network training of the same initial-value problem. The state error, ODE residual, and causal-weight residuals are also checked. The comparison is descriptive because multiple configuration options are changing simultaneously.



\section{Problem and Related Work}
In this section we introduce the mathematical goal of our study and place our configuration within the framework of existing neural-network techniques. We first describe the exact equations and the initial value problem of the Lorenz trajectory that we want the model to learn. Then we explain the existing methods that we used and then we merged them.

\subsection{Lorenz System Initial-Value Problem}
Our main object is a three-state nonlinear ordinary differential equation (ODE) that we need to solve. The Lorenz-1960 system general form \cite{matthews2024pinnde} is:
\begin{equation}
\begin{aligned}
\frac{dx}{dt} &= kl\left(\frac{1}{k^2 + l^2} - \frac{1}{k^2}\right) yz,\\
\frac{dy}{dt} &= kl\left(\frac{1}{l^2} - \frac{1}{k^2 + l^2}\right) xz,\\
\frac{dz}{dt} &= \frac{kl}{2}\left(\frac{1}{k^2} - \frac{1}{l^2}\right) xy.
\end{aligned}
\label{eq:lorenz1960}
\end{equation}
where $k$ and $l$ are the wavenumbers of the two Fourier modes retained by the maximum-simplification truncation. Suppose the state vector is given by $u=(x,y,z)^\mathsf{T}$. Substituting $k=2$ and $l=1$ gives exactly the coefficient triple $(-0.1,\,1.6,\,-0.75)$. The system implemented is:
\begin{equation}
\dot{x}=-0.1\,yz,\qquad \dot{y}=1.6\,xz,\qquad \dot{z}=-0.75\,xy,
\label{eq:ode}
\end{equation}
For our experiment we used the coefficients $(a_x,a_y,a_z)=(-0.1,1.6,-0.75)$. We track the system starting from an initial state \cite{matthews2024pinnde} of $u(0)=(0.5,0.75,1.0)^\mathsf{T}$ until a final time of $T=13.26446$. On this particular rounded endpoint, the reference trajectory nearly returns to its starting state. It is very similar but not exactly closed so we call this interval a nearly closed orbit and not an exact period.


\subsection{Related Work}
In this study we are making an experimental investigation of a series of well known concepts. The approach adopted is a PINN that embeds the physics residual into the training objective \cite{raissi2019}. We use a hard trial-function design \cite{lagaris1998} to ensure our starting state precisely at the point we want to implement it . Learning long trajectories is a well-known failure mode of PINN training, diagnosed from the loss-landscape and neural-tangent-kernel perspectives~\cite{krishnapriyan2021,wang2022ntk,rathore2024}. We take inspiration from the causal residual weighting technique~\cite{wang2024causal}, its causal-sweeping generalization~\cite{penwarden2023}, and gradient-flow remedies~\cite{wang2021gradient}. Benchmark studies further report that performance claims often do not survive strong baselines~\cite{hao2024pinnacle,mcgreivy2024}, which motivates our config-locked, artifact-backed reporting. Lastly segmenting the timeline is based on ideas of causal sweeping and temporal decomposition \cite{penwarden2023}. Our implementation is an investigation of the performance of the individual components in this particular closed orbit of the Lorenz-1960 system.

\section{Windowed PINN Method and Implementation}
This section explains the implementation of windowed PINN with causal weights. This section also contains the model construction, physics enforcement, window segmentation, its handoff and the sequence in how the optimizer works for this particular case. 
\subsection{Model Construction and Time Routing}
We used several modelling strategies to approximate the solution for the Lorenz-1960 system trajectory over  $[0,T]$, where $T=13.26446$ is one closed loop for the chosen initial state. We first tried one network to learn the entire trajectory at once. As it stalled on a plateau, we split the time interval into equal segments.

\begin{figure}[t]
\centering
\includegraphics[width=\columnwidth]{figures/plateau_training.pdf}
\caption{Single-network training on the full interval: the ODE-residual loss plateaus near $2.8\times10^{-4}$ (shaded) over the final ${\sim}6{,}000$ epochs, so late-time error remains large.}
\label{fig:plateau}
\end{figure}

We tried to keep the split window closer to 0.5, as small windows are better \cite{wang2024causal}. In a first experiment, we trained one network sequentially over the 27 segments. Whenever training moved to a new segment, the fit of the earlier segments degraded, which is the handoff-forgetting failure quantified in Section~\ref{sec:comparison}. To address these issues, we introduced WPINN (Windowed PINN). We split the total time $[0,T]$ equally into 27 windows and gave each window its own network. Each window is approximately $0.491$ time units wide. Shorter window gives each network an easier piece to learn, therefore we tried to keep the window size closer to 0.5. 

Each window has its own network  $N_k(t;\theta_k)$ with 1 input layer (t), 4 hidden layers of 60 neurons with tanh unit, and with an output layer with 3 output $(x,y,z)$. For the model parameter initialization we used xavier-uniform weights and zero biases. The 27 networks do not share parameters, but consecutive windows share a boundary. The completed windows predicted endpoint for the boundary is used for the initial state of the next window. 

\subsection{Hard Initial State and Window Handoff}
As the Lorenz-1960 system is an initial value problem, we enforce the initial state of each window through a trial solution.

\begin{equation}
\hat{u}_k(t)
  = u_{0,k} + (t-t_k)N_k(t;\theta_k),
\label{eq:trial}
\end{equation}

where $k=0,\ldots,26$, $t_k$ is the starting time of window $k$, and $u_{0,k}$
is its fixed starting state. At $t=t_k$, the second term is zero, so
$\hat{u}_k(t_k)=u_{0,k}$ regardless of the network parameters
\cite{lagaris1998}. The first window uses the trial initial state; every later window uses the predicted endpoint of the previous window as its initial state.

\begin{equation}
u_{0,k+1}=\hat{u}_k(t_{k+1}), \qquad k=0,\ldots,25.
\label{eq:handoff}
\end{equation}
 
Here window k+1 starts from a copy of the weights of window k then train on its own parameters. This makes the predicted values continuous at the window boundaries by construction. As the starting point of the later windows is predicted, it can carry error from earlier windows. 

\begin{figure}[t]
\centering
\includegraphics[width=\columnwidth]{figures/method.pdf}
\caption{Complete implementation pipeline}
\label{fig:pipeline}
\end{figure}

\subsection{Collocation, Residual, and Causal Loss}
To predict and train the networks of each window, we split the time range into 1536 uniform collocation points per window. Ideally 1535 points for each window and 1 handoff point that becomes the initial state for the next window. While training, all the points in a single window are passed to the network at once which gives an output of shape $(1536,3)$.
Then we apply the trial solution to enforce the intial value. While calculating ODE residuals, pytorch automatic differentiation uses the vectors of x,y,z and create a graph that is differentiable. Pytorch autograd function gives the rate of change for all the collocation points with respect to time. And we use this value and the mathematical derivative to calculate the ODE residual.

\begin{equation}
r_k(t)
  = \frac{d\hat{u}_k(t)}{dt} - f\bigl(\hat{u}_k(t)\bigr).
\label{eq:residual}
\end{equation}

For each ordered collocation time $t_{k,i}$, we take the mean of the three squared residual components:

\begin{equation}
L_{k,i}
  = \frac{1}{3}\left(
      r_{k,x}(t_{k,i})^2+
      r_{k,y}(t_{k,i})^2+
      r_{k,z}(t_{k,i})^2
    \right).
\label{eq:pointloss}
\end{equation}

To solve a long horizon time dependent ODE system, only residual loss is not sufficient. That is why we use causal weights and a causal weighted loss. The causal weight at a time depends on the accumulated residual loss at earlier times \cite{wang2024causal}:

\begin{align}
w_{k,i}
  &= \exp\left(-\epsilon\sum_{j<i}L_{k,j}\right),
     \label{eq:weights}\\
\mathcal{L}_{k,\epsilon}
  &= \frac{1}{n_k}\sum_{i=1}^{n_k}w_{k,i}L_{k,i}.
     \label{eq:causalloss}
\end{align}

Here $n_k=1536$. $L_{k,j}$ in \eqref{eq:weights} means the loss at earlier points than $i$, where $j < i$ and $L_{k,i}$ means the loss at collocation point $i$ in \eqref{eq:causalloss}. $w_{k,i}$ denotes the causal weights of the collocation points. Here the causal weights are fixed while Adam optimizer calculated the gradient. We use four $\epsilon$ stages: $0.01$, $0.1$, $1$, and $10$. A small $\epsilon$ leaves the weights nearly uniform; each stage increases the penalty on late-time residuals until the earlier residual has been reduced, at which point the weights rise toward one.

For a long-horizon system like Lorenz-1960, the residual loss alone is not enough. In the windowed algorithm in particular, early times must be fitted before late times: errors accumulate through the handoff and backtracking is impossible. Later states depend on earlier states, so reducing the earlier error is a precondition for reducing the later error. To address this, we first calculate the per-collocation point loss, $L_{k,i}$ which is the mean of the three squared ode residuals at time $i$. We assign weights for each point in time. Each points weight updates its value with the accumulated loss of all earlier points. The Causal loss \eqref{eq:causalloss} is the weighted mean of the residual loss. Weights are recomputed at every epoch. Later section explains the stopping rules of the stages and unweighted L-BFGS refinement steps. 

\subsection{Optimizer Sequence and Saved State}

The causal loss defined above is minimized separately in each window. Figure~\ref{fig:pipeline}
summarizes the sequence across the windows. Algorithm \ref{alg:training} gives the training steps. Each window creates an Adam optimizer and a StepLR scheduler. The learning rate starts at $10^{-3}$ and is multiplied by 0.9 every 1,000 updates. It is maintained for the four stages of $\epsilon=(0.01,0.1,1,10)$ across the window. There is a 4,000 update limit for each stage. The causal weights are recomputed at every Adam update. If the minimum causal weight over the window's 1,535 interior points exceeds 0.99, training advances to the next $\epsilon$ stage; if the 4,000-step cap is reached first, it advances anyway and records the stage as capped. When all four stages are completed then L-BFGS tries to minimize the loss using unweighted residual mean up to 1,000 epoch. To keep track of training we use a snapshot at every 50 epoch to keep training history.

The code saves the parameters of each window, transfers the endpoint state, and copies the parameters to the next trainable window. Each run writes a config record; rerunning with a different configuration into the same run directory is rejected. The run also outputs stage status and residual snapshots. DOP853 values are captured at set up evaluative points. They do not enter the loss, gradient, or stopping test. After training, a dedicated custom function collects all the artifacts and evaluates the final piecewise model and its residual on the collocation and evaluation grids. Finally, a summary function computes error and invariant measures from these arrays. The reported full residual is one of the evaluation.

\begin{algorithm}[t]
\caption{Sequential causal window training}
\label{alg:training}
\begin{algorithmic}[1]
\Require Ordered windows $W_0,\ldots,W_{26}$ and their uniform time points
\State Set $u_{0,0}\gets (0.5,0.75,1.0)^\mathsf{T}$
\For{$k=0,\ldots,26$}
  \If{$k>0$}
    \State Set $u_{0,k}\gets\hat{u}_{k-1}(t_k)$
    \State Initialize $\theta_k\gets\theta_{k-1}$
  \EndIf
  \State Add $t_{k+1}$ to the residual points if $k<26$
  \For{$\epsilon\in(0.01,0.1,1,10)$}
    \For{at most 4,000 Adam steps}
      \State Network input: collocation time $t_{k,i}\in\mathbb{R}$, 
        \State Here  $i=0,\ldots,1535$
      \State Output: $N_k(t_{k,i};\theta_k)=(x_{k,i},y_{k,i},z_{k,i})\in\mathbb{R}^3$
      \State Compute trial solution: $u_k$ by \eqref{eq:trial}
      \State Compute $r$ by autodiff; sort times and $L_i$
      \State Set $w_i$ by \eqref{eq:weights}; detach $w$
      \State Backpropagate \eqref{eq:causalloss}; Adam then StepLR
      \If{verified post-update min$\_w_{k,i}>0.99$}
        \State \textbf{break}
      \EndIf
    \EndFor
    \State Record whether the threshold was met or capped
  \EndFor
  \State Run L-BFGS on the unweighted residual loss
  \State Save window model, endpoint state, and config record
\EndFor
\end{algorithmic}
\end{algorithm}

\section{Experimental Protocol}
This section explains the experimental setup for the saved windowed run. It contains the configuration and reference information, the evaluation measures and the separate full interval comparison.




\subsection{Baseline Reference}




Table \ref{tab:config} shows the parameters used for the reported windowed run. We use SciPy’s DOP853 \cite{virtanen2020} as the numerical reference, a relative tolerance of $10^{-10}$ and a absolute tolerance of $10^{-12}$. It is sampled at 13,265 equally spaced time points, disjoint from the training grid. The reference enters evaluation only, never the loss, gradient, or stopping test.


\begin{table}[t]
\caption{windowed-run configuration}
\label{tab:config}
\centering\footnotesize
\begin{tabular}{ll}
\toprule
Setting & Value\\
\midrule
Interval; windows & $[0,13.26446]$; 27 equal windows\\
Network per window & $1$--$60$--$60$--$60$--$60$--$3$, tanh\\
Parameters & 11,283 per window\\
Initial condition; precision & 0.5, 0.75, 1.0; float64\\
Uniform collocation grid & 41,446 unique points\\
Adam $\epsilon$ stages & $0.01,0.1,1,10$\\
Stage threshold; epoch limit & min$\_w_{k,i}>0.99$; 4,000\\
Adam rate; scheduler & $10^{-3}$; $\times0.9$ per 1,000 steps\\
L-BFGS budget & at most 1,000 iterations/window\\
Warm start; seed & Initialize $\theta_k\gets\theta_{k-1}$; 0\\
Reference grid & 13,265 uniform times\\
\bottomrule
\end{tabular}
\end{table}

\subsection{Measures and Comparison Run}
On the reference grid, let us define $e(t_j)=\hat{u}(t_j)-u_{\mathrm{ref}}(t_j)$. The mean absolute error, RMSE, and maximum absolute error are reported. The combined RMSE is
\begin{equation}
\operatorname{RMSE}_{2}=\left(\frac{1}{n_{\rm eval}}\sum_j\|e(t_j)\|_2^2\right)^{1/2},
\label{eq:rmse}
\end{equation}
while the combined maximum error is $\max_j\|e(t_j)\|_2$. 
The dense-grid residual mean square is the average of the three residual components squared, taken over the times of the evaluation. The comparison run is with the same network depth and width for the entire interval [0, T]. The windowed run simultaneously changes the number of networks, the grid, loss, warm start and
the optimizer schedule. We measure their final output and compare the two runs.






\section{Results}
This section reports the measurements taken from the completed training artifacts. We first look at what the windowed model does correctly and how well it follows the true path and physics. We then compare these results with the baseline reference and a single-network run for comparison with the behaviour of the model.

\subsection{State Accuracy and Physics Checks}
We compare the windowed PINN with an extremely accurate numerical reference which is SciPy's DOP853. As shown in Table \ref{tab:errors}, our model achieved a combined RMSE of $6.97\times10^{-5}$ across the 13,265 uniform times in our evaluation grid. The highest state error norm was $3.09\times10^{-4}$. Figure \ref{fig:trajectory} plots these predictions against the reference over time. The trajectories overlap throughout, but the error shows small peaks, particularly near $t\simeq 5$.

\begin{table}[t]
\caption{Windowed-run error against DOP853 on 13,265 times}
\label{tab:errors}
\centering\footnotesize
\begin{tabular}{lccc}
\toprule
State & MAE & RMSE & Maximum\\
\midrule
$x$ & $1.89\times10^{-5}$ & $2.74\times10^{-5}$ & $1.13\times10^{-4}$\\
$y$ & $2.94\times10^{-5}$ & $4.96\times10^{-5}$ & $2.88\times10^{-4}$\\
$z$ & $2.78\times10^{-5}$ & $4.05\times10^{-5}$ & $1.30\times10^{-4}$\\
State norm & $5.12\times10^{-5}$ & $6.97\times10^{-5}$ & $3.09\times10^{-4}$\\
\bottomrule
\end{tabular}
\end{table}

\begin{figure}[t]
\centering
\includegraphics[width=\columnwidth]{figures/solution_vs_reference.pdf}
\caption{PINN solution vs reference with signed errors.}
\label{fig:trajectory}
\end{figure}

We also measured how closely the model obeyed the underlying equations. The MSE of the ODE residual was very small on the evaluation grid at $1.872\times10^{-7}$. Also the residual MSE of windowed run on collocation grid is nearly same ($1.874\times10^{-7}$). The DOP853 reference drifted less than $4\times10^{-10}$ for perspective.

Lastly, the hard trial solution ensures that the values of the state are exactly continuous at the inner 26 interfaces connecting the windows.


\subsection{Comparison and Training Record}
\label{sec:comparison}
The differences with respect to the two saved configurations are shown in Table \ref{tab:comparison}. The RMSE value of the entire interval network is 2.237, which is much larger than the windowed run network's RMSE value of $6.97\times10^{-5}$. This gap is the central observation for the recorded settings of this initial-value problem. All of the above are uncontrolled types of ablation in the concept of causal weighting \cite{wang2024causal}, windowing \cite{penwarden2023} and warm starts as well as precision and even L-BFGS. In particular the model with the window is much more complicated with 27 times as many parameters and a different collocation grid. Reducing the parameter count of each network might result in similar error.

\begin{table}[t]
\caption{Saved full-interval and windowed runs on the same IVP}
\label{tab:comparison}
\centering\footnotesize
\begin{tabular}{lcc}
\toprule
Measure & One network & 27 windows\\
\midrule
Trainable parameters & 11,283 & 304,641\\
Training points & 39,793 LHS & 41,446 uniform\\
Adam budget & 40,000 & 4 stages $\times\,\le$4,000\\
L-BFGS iterations & 5,000 & 1,000\\
LR schedule & $\times0.9$ per 5,000 & $\times0.9$ per 1,000\\
Combined state RMSE & 2.237 & $6.97\times10^{-5}$\\
Maximum state-error norm & 3.346 & $3.09\times10^{-4}$\\
Collocation residual MSE & $2.80\times10^{-4}$ & $1.87\times10^{-7}$\\
Recorded wall time (s) & 4,529 & 881\\
\bottomrule
\end{tabular}
\end{table}

The 27 run windows had 34,765 Adam updates and 1000 L-BFGS closure evaluations per window. Closure evaluations are function calls made by L-BFGS, not iterations. Only two of the 108 causal stages ended at the 4,000-step cap without meeting the $0.99$ weight threshold.
\section{Discussion and Limitations}
There are a few limitations of our method. The trial solution \cite{lagaris1998} defines the initial state for each window exactly. The values remain continuous at the boundaries, but the derivatives can still change suddenly at the boundaries. This comparison with the single network is only a description of two saved runs. The two runs differed in a number of things, including the number of networks, grid, loss and schedule. The original comparison did not include a windowed run without causal weights, because \cite{wang2024causal} report one to two orders of magnitude lower error when causality is respected; a budget-matched ablation isolating the causal weights accompanies the released run artifacts. The reported windowed run uses a single seed, one initial state, and one coefficient set; seed sensitivity is quantified in the released run artifacts. Other initial states and other coefficient sets may result in different behavior. There is a weakness with the handoff. All windows that follow the first window begin from a predicted state and not the actual state. This means that errors from earlier windows get carried over to later windows. Only 2 of the 108 causal stages were at the 4,000-step limit but failed to reach the $0.99$ weight limit. The causal weights Wang \emph{et al.} vary depending on the number and spacing of the collocation points so if we increase the number of collocation points that will decrease the training loss. To obtain a fairer future test, one and only one training option will be changed at a time, with several seeds.

\section{Conclusion}



We trained a 27-window causal PINN over one nearly closed orbit of the Lorenz-1960 three-wave system, $t\in[0,13.26446]$, with hard initial-state enforcement and sequential window handoff. The saved run reaches a combined state RMSE of $6.97\times10^{-5}$ against DOP853, while a single network on the full interval stalls at a residual plateau (RMSE 2.237). The reference solution is verified independently by conserved-quadratic and cross-solver checks, and every number reported here is traceable to a config-locked run artifact. Future work isolates each training option one at a time under a multi-seed protocol and generalizes to problems with boundary conditions at both ends.








\bibliographystyle{IEEEtran}
\bibliography{references}

\end{document}
